\icode[magenta]{N25E-EME-MU-L4-260625-002}
\nceeID[4]{2025湖北}{10}
如图所示，在\(xOy\)平面内有一以\(O\)点为中心的正五边形，顶点到\(O\)点的距离为\(R\)。在正五边形的顶点上顺时针方向依次固定电荷量为\(q\)、\(2q\)、\(3q\) 、\(4q\)、\(5q\)的正点电荷，且电荷量为\(3q\)的电荷在\(y\)轴正半轴上。静电力常量为\(k\)，则\(O\)点处的电场强度\paren[]

\begin{tikzpicture}[font=\scriptsize,>=stealth]
	\draw[->,blue] (-2,0)--(2,0)node[right]{\(x\)};
	\draw[->,blue] (0,-2)--(0,2)node[above]{\(y\)};
	\node at(.2,-.2)[blue]{\(O\)};
	\draw (18:1.2)node[right]{\(4q\)}--(90:1.2)node[above left]{\(3q\)}--(162:1.2)node[left]{\(2q\)}--(234:1.2)node[left]{\(q\)}--(306:1.2)node[right]{\(5q\)}--cycle;
\end{tikzpicture}
\begin{choices}
	\item 方向沿\(x\)轴负方向
	\item 方向与\(x\)轴负方向成\qty{18}{\degree}夹角斜向下
	\item 大小为\(\frac{2kq}{R^2}(\cos\text{\,\qty{54}{\degree}}+\cos\text{\,\qty{18}{\degree}})\)
	\item 大小为\(\frac{2kq}{R^2}(2\cos\text{\,\qty{54}{\degree}} +\cos\text{\,\qty{18}{\degree}})\)
\end{choices}
\begin{answer}
	AD
\end{answer}